% ============================================================================
% SEC02.TEX - Section 11.2: Membuat Enemy Ship
% ============================================================================

\section{Membuat Enemy Ship}

\subsection{Setup Prefab}
\begin{enumerate}
    \item Buat Sprite baru (Misalnya bentuk V atau pesawat merah). Beri nama \texttt{EnemyShip}.
    \item Tambahkan \texttt{Rigidbody 2D} (Gravity 0) dan \texttt{Polygon Collider 2D}.
    \item \textbf{PENTING}: Set Tag menjadi \texttt{Enemy} (agar bisa menghancurkan player dan dihancurkan peluru).
    \item Jadikan Prefab.
\end{enumerate}

\begin{figure}[ht]
    \centering
    \begin{tikzpicture}[node distance=2.5cm, auto]
        \node (spawn) [draw, rounded corners] {Spawn};
        \node (chase) [draw, rectangle, right=of spawn] {Chase Player};
        \node (rotate) [draw, rectangle, right=of chase] {Rotate to Target};
        
        \draw[->] (spawn) -- (chase);
        \draw[->, loop above] (chase) to node {Update Position} (chase);
        \draw[->] (chase) -- (rotate);
    \end{tikzpicture}
    \caption{Logika AI Sederhana}
    \label{fig:enemy_ai}
\end{figure}
